% Final: full proof and exact-primary-source audits passed by root and frobenius.
\paragraph{C.3ad (Proved in the stated rational smooth-cone scope): the top quotient over algebraic numbers.}
Let $k=\overline{\mathbb Q}$, let $V/k$ be a smooth projective
integral variety of dimension $d\ge2$, and let $L$ be very ample.
Assume its full section ring is normal and generated in degree one.
Let $R$ be the strict henselization of the local ring of its affine
cone at the vertex, and put
\[
             X=\operatorname{Spec}R\setminus\{\mathfrak m_R\}.
\]
Then the natural maps give
\begin{equation}\label{eq:C3ad-top-cone}
 \mathrm{CH}^d(X)_{\mathbb Q}\simeq
 \frac{\mathrm{CH}_0(V)_{\mathbb Q}}
      {c_1(L)\mathrm{CH}^{d-1}(V)_{\mathbb Q}},
 \qquad B^d(X)_{\mathbb Q}=0.
\end{equation}
The equivalence relation defining $B^d$ remains ordinary
algebraic equivalence on actual fixed products. In dimensions
$d>2$ this statement does not assert integral finiteness;
it also does not concern arbitrary strict local singularities.

\begin{proof}
Use the resolution $Z\to\operatorname{Spec}R$, its exceptional
divisor $i:V\hookrightarrow Z$, and the projection $p:Z\to V$
from C.3e and C.3ab. Thus $Z$ is regular of dimension $d+1$,
$p i=1$, $Z\setminus V=X$, and
\[
 V_a=Z\times_R R/\mathfrak m_R^a
     =\operatorname{Spec}_V
           (\mathcal O_V\oplus L\oplus\cdots\oplus L^{a-1}).
\]
The multiplication on the right is truncated graded
multiplication. The pro cdh argument in C.3ab applies without
change in this dimension: $K_0(R)=K_0(R/\mathfrak m_R^a)=\mathbb Z$
and the first map is the graph of $[\mathcal O_Z]$.
Consequently
\begin{equation}\label{eq:C3ad-finite-detection}
               K_0(Z)\longrightarrow K_0(V_a)
               \quad\text{is injective for }a\gg0.
\end{equation}

\emph{The exact cyclic-homology input.}
For an arbitrary commutative $\mathbb Q$-algebra $A$, including
a nonregular algebra, the Eulerian decomposition satisfies
\begin{equation}\label{eq:C3ad-Hodge}
 HC_n^{(r)}(A)=0\quad(r>n\ge0),\qquad
 HC_n^{(n)}(A)=\Omega^n_{A/\mathbb Q}/d\Omega^{n-1}_{A/\mathbb Q}.
\end{equation}
These are \citet[Theorems 4.6.7--4.6.8,
pp.~151--152]{D_Loday1992}. The hypotheses are commutativity
and a ground ring containing $\mathbb Q$; no smoothness
assumption occurs. The displayed top component is obtained
by taking the cokernel of the Hochschild boundary and then
the Connes operator, as in the source's proof.
Thus \eqref{eq:C3ad-Hodge} does not apply the smooth HKR
isomorphism to $V_a$.

The relative Chern character preserves the Adams decomposition,
with the shift
\[
       K_q^{(r)}(A,I)_{\mathbb Q}
                    \simeq HC_{q-1}^{(r-1)}(A,I)
\]
for a nilpotent ideal $I$. We use the scheme-compatible
form of Cathelineau's theorem proved by
\citet[Theorem 6.1 and Corollary 6.2,
author p.~15]{D_CortinasHaesemeyerWeibel2009}.
The shift is explicitly also the one in their Theorem B.10.
The relative groups here are already uniquely divisible,
but we retain rational coefficients throughout this proof.
Since the projection $V_a\to V$ splits restriction,
the relative groups are kernels of the absolute restriction
maps. Equations \eqref{eq:C3ad-Hodge} therefore imply, for
the relative Zariski sheaves on $V$,
\begin{equation}\label{eq:C3ad-relative-sheaves}
 \mathcal K_q^{(d)}(V_a,V)_{\mathbb Q}=0\quad(0\le q<d),
 \qquad
 \mathcal K_d^{(d)}(V_a,V)_{\mathbb Q}
       \simeq\bigoplus_{j=1}^{a-1}
                         \Omega^{d-1}_{V/\mathbb Q}\otimes L^j.
\end{equation}
For $q=0$ use nilpotent invariance of $K_0$ for affine rings.

Here is the local calculation and its gluing in the last
isomorphism. Put $s=d-1$ and trivialize $L$, so the algebra
is $A[t]/t^a$. Its positive internal-degree-$j$ forms, for
$1\le j<a$, have the form
\[
              t^j\omega+t^{j-1}dt\wedge\eta,
 \qquad \omega\in\Omega_A^s,
        \quad\eta\in\Omega_A^{s-1}.
\]
Modulo exact forms this has the unique horizontal
representative
\begin{equation}\label{eq:C3ad-horizontal}
              t^j(\omega-j^{-1}d\eta).
\end{equation}
Indeed subtract $j^{-1}d(t^j\eta)$; the remaining exact
relations from a form containing $dt$ map to zero because
$d^2=0$. The relation $a t^{a-1}dt=0$ removes internal
degree $a$, and there are no higher degrees. Inclusion
of horizontal forms is consequently an isomorphism onto
this quotient. Under $t=u t'$ a horizontal $s$-form
transforms as $t^j\omega=u^j(t')^j\omega$.
The isomorphism therefore glues with precisely $L^j$,
without any additional conormal term.

\emph{The weight-$d$ part on the resolution.}
The relative descent spectral sequence, with the Adams
component just specified, is
\[
 H^b(V,\mathcal K_q^{(d)}(V_a,V)_{\mathbb Q})
                         \Longrightarrow K_{q-b}^{(d)}(V_a,V)_{\mathbb Q}.
\]
It follows either by taking the relative fibre of
\citet[Theorem 10.3]{D_ThomasonTrobaugh1990} and the compatible
Adams summand, or directly from the scheme comparison above.
Since $\dim V=d$, only $0\le b\le d$ occur. On the
diagonal $q=b$ all terms vanish by
\eqref{eq:C3ad-relative-sheaves}, except possibly
\[
 H^d\!\left(V,
      \bigoplus_{j=1}^{a-1}\Omega^{d-1}_{V/\mathbb Q}\otimes L^j\right).
\]
The latter vanishes as well. Indeed
$\Omega^1_{k/\mathbb Q}=0$, so these are relative
differentials over $k$, and Kodaira--Akizuki--Nakano
vanishing applies because
\[
                     (d-1)+d>d\qquad(d\ge2).
\]
The exact characteristic-zero version is
\citet[Theorem 1.3(a) and Section 8]{D_EsnaultViehweg1992}.
It follows that
\[
                     K_0^{(d)}(V_a,V)_{\mathbb Q}=0.
\]
The split projection to $V$, together with
\eqref{eq:C3ad-finite-detection}, now proves
\begin{equation}\label{eq:C3ad-weight-isomorphism}
              i^*:K_0^{(d)}(Z)_{\mathbb Q}
                   \xrightarrow{\sim}K_0^{(d)}(V)_{\mathbb Q},
                        \qquad(i^*)^{-1}=p^*.
\end{equation}

\emph{Chow groups and ordinary algebraic equivalence.}
Rational Chern character identifies the weight-$d$
summands in \eqref{eq:C3ad-weight-isomorphism} with
$\mathrm{CH}^d(Z)_{\mathbb Q}$ and
$\mathrm{CH}_0(V)_{\mathbb Q}$, respectively. Its application
to this $Z$ is justified by the same filtered pointed-etale
models used in C.3ab: each model is a smooth variety, and
perfect complexes, cycles, rational-equivalence data and
Chern classes descend to a finite stage. The usual
rational Chern character isomorphism on those models
\citep[Example 15.2.16, pp.~294--295]{frC4Fulton1998} therefore passes to
the limit and remains compatible with $i^*$ and $p^*$.

Chow localization for $V\subset Z$, and
$i^!i_*=c_1(L^{-1})$, give the first isomorphism
in \eqref{eq:C3ad-top-cone}. The maps are the same actual
flat pullbacks as in C.3e. Thus they preserve ordinary
algebraic equivalence, and C.3e also gives injectivity
on the corresponding ordinary quotients. In the top
degree on the connected smooth projective variety $V$,
ordinary equivalence is degree, so $B^d(V)=\mathbb Z$.
The degree of $c_1(L)\cdot C$ is positive for every
integral curve $C\subset V$. Such a curve exists by
hyperplane sections. Hence the image of
$c_1(L)B^{d-1}(V)$ is nonzero in $\mathbb Z$ and is
all of $\mathbb Q$ after tensoring with $\mathbb Q$.
This proves $B^d(X)_{\mathbb Q}=0$.

The last use of rational Chern character is precisely
why the argument, as stated, does not prove integral
finiteness for $d>2$. For $d=2$, the independent
codimension-two integral calculation in C.3ab gives
the stronger finite cyclic answer.
\end{proof}
